\section{Abstract-Stage Triage}
\label{sec:ECTriage}

The venue sweep returns \(379\) records at the eleven venues. Most are not about contextual
optimization at all: the query is deliberately recall-oriented (\cref{sec:ECSearch}), and terms such
as \emph{task loss} and \emph{differentiable optimization} are common in unrelated machine-learning
work. Before obtaining full texts we therefore apply a single coarse filter, which we call
\emph{abstract-stage triage}.

We stress at the outset what this stage does \emph{not} do. It does not apply the inclusion rules of
\cref{sec:ECInclusion}, and it does not determine whether a paper studies a stochastic parameter. It
decides only which full texts are worth obtaining. Every substantive classification in this audit is
made from the full text, for reasons documented below.

\subsection{Why the Filter Is Not a Keyword Rule}

The natural implementation is a keyword or topic rule applied to the abstract, and we initially
adopted one: exclude a paper when its abstract indicates that no parameter is predicted, so that the
optimization instance is fully observed and the learning task is to solve it faster. Papers of this
kind -- amortized optimization, learned warm starts, neural combinatorial solvers -- are genuinely
out of scope, since with no unobserved \(Y\) there are no two oracles and no premium.

That rule fails systematically, and it fails on exactly the papers we can least afford to lose. We
tested it against twelve papers already known to be in scope, judging each from its abstract alone.
Four of the twelve would have been excluded: \citet{berthet2020learning},
\citet{poganvcic2020differentiation}, \citet{sahoo2023backpropagation} and
\citet{niepert2021implicit}. Each describes its
\emph{method} -- differentiating through a combinatorial solver -- and leaves the \emph{setting}, in
which a cost vector is predicted from context, to its experiments. Their abstracts contain no
language about prediction at all.

Abstracts at machine-learning venues foreground methodological contribution, so the papers most
likely to be misread are precisely those in the differentiable-optimization lineage that the search
was widened to capture. A referee has no way to
detect such an exclusion from the published record, since an excluded paper leaves no trace. We
therefore adopt a rule that never asks whether anything is predicted.

\subsection{The Triage Rule}

A record is retained when both of the following hold: an optimization or decision problem is the
object of the work, and a learned component appears in the same pipeline. A record is excluded when
there is no decision problem at all, when the only optimization is machinery for training or
compressing a model (hyperparameter search, architecture search, pruning, distillation), or when the
work is numerical simulation of a physical system. Records that cannot be judged are retained.

The asymmetry is deliberate. A record wrongly retained costs one full-text retrieval; a record
wrongly excluded is invisible for the remainder of the study.

\subsection{Implementation and Reproducibility}

The rule requires reading an abstract for what a paper is about, which no lexical pattern captures.
We implement it as a language-model classifier under a fixed prompt, and treat the prompt as part of
the method rather than as an implementation detail: it is versioned in the replication repository,
states the criterion as enumerated inclusions and exclusions rather than as an abstract principle,
and carries five worked examples with their reasoning, among them \citet{berthet2020learning} as a
retained record. Two instructions in the prompt exist solely to prevent the failure above, and each
records the evidence that motivated it.

The model identifier is pinned and decoding is greedy. For each record we store the assigned
bucket, a one-sentence rationale in the model's own words, a self-reported confidence, the model
identifier, a hash of the prompt, and the run date. The rationale is what makes the stage auditable:
a reader disagreeing with a particular exclusion can see the stated ground for it.

We are careful not to overclaim here. Greedy decoding reduces run-to-run variation but does not
eliminate it, and a provider-side model revision can move a borderline record. This is why the stage
is validated on every execution against a labelled control set of papers known to be in scope,
four of which -- the four named above -- carry no prediction language in their abstracts and are the
binding test; the pipeline halts if any control record is excluded. We report the control-set result
alongside the triage counts.

\subsection{Scope of the Claim}

Triage is a retrieval decision and we make no stronger claim for it. Its errors in the retaining
direction are harmless, being corrected at full text; its errors in the excluding direction are the
ones that matter, and the control set bounds them only on a set of twelve papers, which is small.
Perhaps the most honest statement is that this stage is designed so that its failures are
recoverable, not so that it is correct: every subsequent stage of the audit works from full texts,
and the screening log records a reason for each exclusion so that a reader may disagree with any of
them individually.
